
\subsection{同伦等价与算子 $\mathfrak q$ II：辛微分同胚下的不变性 1}
\label{subsec:4-6-5}
\liftIndex{I180}

本小节与下一小节证明定理 4.6.52。
设 $(M,\omega)$、$(M',\omega')$ 为辛流形，
$\psi\colon M\to M'$ 为辛微分同胚。
令 $L\subset M$ 为具有相对 spin 结构的拉格朗日子流形，
$L'=\psi(L)$。
在 $M,L$ 与 $M',L'$ 上固定相容的几乎复结构、
扰动以及可数的光滑奇异单形集等各种选择，
从而定义
$(C(L;\Lambda_{0,\mathrm{nov}}),\mathfrak m)$
和
$(C(L';\Lambda_{0,\mathrm{nov}}),\mathfrak m)$。
对
$[\mathfrak b]\in H^2(M;\Lambda^+_{0,\mathrm{nov}})$
及
$[\mathfrak b']\in H^2(M';\Lambda^+_{0,\mathrm{nov}})$，
我们已得到滤过 $A_\infty$ 代数
$(C(L;\Lambda_{0,\mathrm{nov}}),\mathfrak m^{\mathfrak b})$
和
$(C(L';\Lambda_{0,\mathrm{nov}}),\mathfrak m^{\mathfrak b'})$；
它们的同伦型分别只依赖于
$\mathfrak b$、$\mathfrak b'$ 的上同调类。

\begin{theorem}{4.6.52}\label{thm:4-6-52}
对每个
$[\mathfrak b]\in H^2(M;\Lambda^+_{0,\mathrm{nov}})$，
存在滤过 $A_\infty$ 代数的同伦等价
\[
\psi_*^{\mathfrak b}\colon
(C(L;\Lambda_{0,\mathrm{nov}}),\mathfrak m^{\mathfrak b})
\longrightarrow
(C(L';\Lambda_{0,\mathrm{nov}}),\mathfrak m^{\psi_*\mathfrak b}).
\]
它在定理 4.1.2 的意义下不依赖于辛同痕的选择。
\end{theorem}

证明定理 4.6.52 之前，先导出它的几个推论。
下面的推论就是引言中定理 B 的 (B.3)。
\begin{corollary}{4.6.53}\label{cor:4-6-53}
如上所述，满足 $L'=\psi(L)$ 的辛微分同胚
$\psi\colon M\to M'$ 诱导映射 $\psi_*$，
使图式 4.6.4 中两个图都交换：
\begin{center}
\centering
\[
\begin{gathered}
\begin{array}{ccc}
\mathcal M_{\mathrm{weak,def}}(L)
 &\xrightarrow{\ \psi_*\ }&
\mathcal M_{\mathrm{weak,def}}(L')\\[-1mm]
\big\downarrow{\scriptstyle\pi_{\mathrm{amb}}}
 &&\big\downarrow{\scriptstyle\pi_{\mathrm{amb}}}\\
H^2(M;\Lambda^+_{0,\mathrm{nov}})
 &\xrightarrow{\ \psi_*\ }&
H^2(M';\Lambda^+_{0,\mathrm{nov}})
\end{array}
\\[3mm]
\begin{array}{ccc}
\mathcal M_{\mathrm{weak,def}}(L)
 &\xrightarrow{\ \psi_*\ }&
\mathcal M_{\mathrm{weak,def}}(L')\\[-1mm]
\big\downarrow{\scriptstyle\mathfrak{PO}}
 &&\big\downarrow{\scriptstyle\mathfrak{PO}}\\
\Lambda_{0,\mathrm{nov}}^{(0)}
 &\xrightarrow{\ \mathrm{id}\ }&
\Lambda_{0,\mathrm{nov}}^{(0)} .
\end{array}
\end{gathered}
\]
\renewcommand{\thefigure}{4.6.4}
\refstepcounter{figure}
{\normalfont\bfseries 图式~\thefigure}\label{diag:4-6-4}
\end{center}
其中的水平箭头均为同构。
图式 4.6.4 左图的下方箭头是
$(\psi^{-1})^*$，即 $\psi^{-1}$ 的拉回。
\end{corollary}

\begin{proof}[推论 4.6.53 的证明]
可写成
\[
\begin{aligned}
\mathcal M_{\mathrm{weak,def}}(L)
=\bigl\{(\mathfrak b,b)\ \bigm|\;&
\mathfrak b\in H^2(M;\Lambda^+_{0,\mathrm{nov}}),\\
&b\in\mathcal M_{\mathrm{weak}}
  (C(L;\Lambda_{0,\mathrm{nov}}),
   \mathfrak m^{\mathfrak b})\bigr\}\big/\sim .
\end{aligned}
\]
由定理 4.6.52，对上述每个 $[\mathfrak b]$
有同伦等价 $\psi_*^{\mathfrak b}$。
根据定理 4.3.22，函子 $\mathcal M_{\mathrm{weak}}$
在同伦下不变，因此
\[
\begin{aligned}
(\psi_*^{\mathfrak b})_*\colon\;
&\mathcal M_{\mathrm{weak}}
 (C(L;\Lambda_{0,\mathrm{nov}}),\mathfrak m^{\mathfrak b})
\\[-1mm]
&\longrightarrow
\mathcal M_{\mathrm{weak}}
 (C(L';\Lambda_{0,\mathrm{nov}}),
  \mathfrak m^{\psi_*\mathfrak b})
\end{aligned}
\]
是双射，且
$(\psi_*^{\mathfrak b})_*(b)
=\psi_*^{\mathfrak b}(e^b)$；
另见引理 3.6.36。
现在定义
\[
\psi_*(\mathfrak b,b)
=\bigl(\psi_*\mathfrak b,
       (\psi_*^{\mathfrak b})_*(b)\bigr).
\]
这里 $\psi_*\mathfrak b$ 是
$\psi^{-1}$ 对 $\mathfrak b$ 的拉回。
由引理 4.3.23 可知，
该映射具备所需性质。
\end{proof}

\begin{proof}[定理 B (B.4) 的证明]
沿用定理 B 的记号。
设 $\boldsymbol b_i=(\mathfrak b,b_i)$
（$i=0,1$）。
由定理 4.6.52 得到同伦等价
\[
\psi_*^{\mathfrak b}\colon
(C(L;\Lambda_{0,\mathrm{nov}}),\mathfrak m^{\mathfrak b})
\longrightarrow
(C(L';\Lambda_{0,\mathrm{nov}}),
 \mathfrak m^{\psi_*\mathfrak b}),
\]
因而也有其在定界余链模空间上诱导的映射
$(\psi_*^{\mathfrak b})_*$。
记
$\psi_*\boldsymbol b_i
=(\psi_*\mathfrak b,(\psi_*^{\mathfrak b})_*(b_i))$；
在涉及边界形变的下列公式中，
简记 $(\psi_*^{\mathfrak b})_*(b_i)$ 为 $\psi_*b_i$。
令
\[
\begin{aligned}
\delta_{b_1,b_0}(x)
 &=\mathfrak m^{\mathfrak b}(e^{b_1}xe^{b_0}),\\
\delta_{\psi_*b_1,\psi_*b_0}(y)
 &=\mathfrak m^{\psi_*\mathfrak b}
   (e^{\psi_*b_1}ye^{\psi_*b_0}).
\end{aligned}
\]
定义复形映射
\[
\begin{aligned}
\psi_*^{\boldsymbol b_1,\boldsymbol b_0}\colon\;
&(C(L;\Lambda_{0,\mathrm{nov}}),\delta_{b_1,b_0})\\
&\longrightarrow
(C(L';\Lambda_{0,\mathrm{nov}}),
 \delta_{\psi_*b_1,\psi_*b_0})
\end{aligned}
\]
为
$\psi_*^{\boldsymbol b_1,\boldsymbol b_0}(x)
=\psi_*^{\mathfrak b}(e^{b_1}xe^{b_0})$。
容易看出它是链映射，故诱导
\begin{equation}
\begin{aligned}
(\psi_*^{\boldsymbol b_1,\boldsymbol b_0})_*\colon\quad
&HF((L,\boldsymbol b_1),(L,\boldsymbol b_0);
     \Lambda_{0,\mathrm{nov}})\\
&\longrightarrow
HF((L',\psi_*\boldsymbol b_1),
   (L',\psi_*\boldsymbol b_0);
   \Lambda_{0,\mathrm{nov}}).
\end{aligned}
\tag{4.6.54}\label{eq:4-6-54}
\end{equation}
注意
$\psi_*^{\boldsymbol b_1,\boldsymbol b_0}
\equiv\psi_*\pmod{\Lambda^+_{0,\mathrm{nov}}}$，
而 $\psi_*$ 在 $\overline{\mathfrak m}_1$
上同调上诱导同构。
命题 5.2.38 (2)（将在第 5.2 节证明）表明
(4.6.54) 是同构。
也可用第 6 章将建立的谱序列证明这一点。
定理 B (B.4) 得证。
\end{proof}

为证明定理 4.6.52，
须结合定理 4.6.1 的证明和第 3.8.6 小节的构造。
先把可数生成子复形等选择说清楚。
如定理 3.8.32，
选取 $L$ 的可数光滑奇异单形集
$\mathcal X_{1+}(L)\supset\mathcal X_1(L)$
以及 $M$ 的单形集 $\mathcal X(M)$，
使每个 $Q\in\mathcal X(M)$ 都与 $L$ 横截，
且适当剖分后的 $Q\cap L$ 的各个单形
属于 $\mathcal X_1(L)\subset\mathcal X_{2+}(L)$。
于是算子
\[
\mathfrak q_{\ell,k}\colon
E_\ell C(M;\Lambda_{0,\mathrm{nov}})[2]
\otimes B_kC_{1+}(L;\Lambda_{0,\mathrm{nov}})[1]
\longrightarrow
C_{1+}(L;\Lambda_{0,\mathrm{nov}})[1]
\]
有定义；其中 $C(M;\Lambda_{0,\mathrm{nov}})$、
$C_{1+}(L;\Lambda_{0,\mathrm{nov}})$
分别由 $\mathcal X(M)$、$\mathcal X_{1+}(L)$ 生成。

同样可在 $M',L'$ 上选择
$\mathcal X(M')$、$\mathcal X_{2+}(L')$，
使
\[
\mathfrak q_{\ell,k}\colon
E_\ell C(M';\Lambda_{0,\mathrm{nov}})[2]
\otimes B_kC_{2+}(L';\Lambda_{0,\mathrm{nov}})[1]
\longrightarrow
C_{2+}(L';\Lambda_{0,\mathrm{nov}})[1]
\]
有定义；两个复形分别由
$\mathcal X(M')$、$\mathcal X_{2+}(L')$ 生成。
\footnote{译者注：原书实体页 272
在此把第二个生成集写为
$\mathcal X_{1+}(L')$，
但它刚指定的是 $\mathcal X_{2+}(L')$，
且 $\mathcal X_{1+}$ 只在 $L$ 上定义。
见疑误 A-0250。}
还可选取包含
$\psi(\mathcal X(M))$ 与 $\mathcal X(M')$
的 $\mathcal X'(M')$，
使其生成的 $C'(M';\Lambda_{0,\mathrm{nov}})$
满足
\[
H^*(C'(M';\Lambda_{0,\mathrm{nov}}))
\cong H^*(M';\Lambda_{0,\mathrm{nov}}).
\]
在这些选择下，有如下命题。

\begin{proposition}{4.6.55}\label{prop:4-6-55}
存在 $L'$ 的可数光滑奇异单形集
$\mathcal X_{3+}(L')$，
它包含定理 4.6.1 中的 $\mathcal X_3(L')$
及上述 $\mathcal X_{2+}(L')$，
并生成 $C_{3+}(L';\Lambda_{0,\mathrm{nov}})$。
还存在次数为 $0$ 的一列算子
\[
\begin{aligned}
\mathfrak f_{\ell,k}\colon\quad
&E_\ell C(M;\Lambda_{0,\mathrm{nov}})[2]
\otimes B_kC_{1+}(L;\Lambda_{0,\mathrm{nov}})[1]\\
&\longrightarrow C_{3+}(L';\Lambda_{0,\mathrm{nov}})[1],
\end{aligned}
\]
具有下列性质。

\noindent\textup{(4.6.56.1)}
$\mathfrak f_{0,k}=\mathfrak f_k$，
其中 $\mathfrak f_k$（$k=0,1,2,\ldots$）
是定义 4.6.12 中滤过 $A_\infty$ 同态的分量。

\noindent\textup{(4.6.56.2)}
原先定义在 $C_{2+}(L')$ 上的算子
$\mathfrak q_{\ell,k}$ 可扩张为
\[
\begin{aligned}
\mathfrak q_{\ell,k}\colon\quad
&E_\ell C'(M';\Lambda_{0,\mathrm{nov}})[2]
\otimes B_kC_{3+}(L';\Lambda_{0,\mathrm{nov}})[1]\\
&\longrightarrow C_{3+}(L';\Lambda_{0,\mathrm{nov}})[1],
\end{aligned}
\]
并满足 (3.8.33)。
\footnote{译者注：原书实体页 272
把本条扩张前后的靶写为
$BC_{2+}[1]$、$BC_{3+}[1]$。
依定理 3.8.32 及本页上方算子
$\mathfrak q_{\ell,k}$ 的定义，
各 $\mathfrak q_{\ell,k}$ 是单输出分量，
靶应为 $C_{2+}[1]$、$C_{3+}[1]$；
余代数层面的扩张另由 (3.8.33) 给出。
见疑误 A-0251。}

\noindent\textup{(4.6.56.3)}
定义
\[
\begin{aligned}
\widehat{\widehat{\mathfrak f}}\colon\quad
&EC(M;\Lambda_{0,\mathrm{nov}})[2]
\otimes BC_{1+}(L;\Lambda_{0,\mathrm{nov}})[1]\\
&\longrightarrow BC_{3+}(L';\Lambda_{0,\mathrm{nov}})[1],
\end{aligned}
\]
使其在
$B_mC_{3+}(L';\Lambda_{0,\mathrm{nov}})[1]$
中的分量为
\[
\begin{aligned}
(\pi_m\circ\widehat{\widehat{\mathfrak f}})
 (\boldsymbol y\otimes\boldsymbol x)
=\sum_{c_1,c_2}(-1)^{\epsilon(c_1,c_2)}
&\mathfrak f(
 \boldsymbol y_{c_1}^{(m;1)}
 \otimes\boldsymbol x_{c_2}^{(m;1)})\\[-1mm]
{}\otimes\cdots\otimes{}&
\mathfrak f(
 \boldsymbol y_{c_1}^{(m;m)}
 \otimes\boldsymbol x_{c_2}^{(m;m)}),
\end{aligned}
\]
其中\footnote{译者注：原书实体页 272
在 $\epsilon(c_1,c_2)$ 中把第二因子的下标写成
$c_1$；但对应的 $\boldsymbol x$ 分解
在前式中由 $c_2$ 标记。
见疑误 A-0252。}
\[
\epsilon(c_1,c_2)
=\sum_{1\le j<i\le m}
\deg\boldsymbol y_{c_1}^{(m;i)}
\deg'\boldsymbol x_{c_2}^{(m;j)} .
\]
那么有\footnote{译者注：原书实体页 272
在 (4.6.57) 第一项求和中写
$\sum_c$，却把同一二重分解的第一因子
标作 $\boldsymbol y_{c_1}^{(2;1)}$。
这里与第二因子统一为下标 $c$；
见疑误 A-0254。}
\begin{equation}
\begin{aligned}
&\mathfrak f(\widehat\delta_M(\boldsymbol y)
             \otimes\boldsymbol x)\\
&\quad+\sum_c
\mathfrak q\bigl(
 \psi_*^\otimes(\boldsymbol y_c^{(2;1)})
 \otimes
 \widehat{\widehat{\mathfrak f}}
 (\boldsymbol y_c^{(2;2)}\otimes\boldsymbol x)
 \bigr)\\
&\quad-\sum_{c_1,c_2}
(-1)^{
 \deg\boldsymbol y_{c_1}^{(2;2)}
 \deg'\boldsymbol x_{c_2}^{(3;1)}
 +\deg'\boldsymbol x_{c_2}^{(3;1)}
 +\deg\boldsymbol y_{c_1}^{(2;1)}
}\\
&\qquad\mathfrak f\Bigl(
 \boldsymbol y_{c_1}^{(2;1)}\otimes
 \bigl(
 \boldsymbol x_{c_2}^{(3;1)}
 \otimes\mathfrak q(
   \boldsymbol y_{c_1}^{(2;2)}
   \otimes\boldsymbol x_{c_2}^{(3;2)})
 \otimes\boldsymbol x_{c_2}^{(3;3)}
 \bigr)\Bigr)=0 .
\end{aligned}
\tag{4.6.57}\label{eq:4-6-57}
\end{equation}
这里
$\psi_*^\otimes\colon
EC(M;\Lambda_{0,\mathrm{nov}})[1]
\to EC'(M';\Lambda_{0,\mathrm{nov}})[1]$
由
$\psi_*\colon C(M;\Lambda_{0,\mathrm{nov}})[1]
\to C'(M';\Lambda_{0,\mathrm{nov}})[1]$
诱导。
\end{proposition}

命题 4.6.55 留待下一小节证明。
这里说明它如何推出定理 4.6.52。
定义
\[
\begin{aligned}
\psi_*^{\mathfrak b}\colon\quad
&BC_{1+}(L;\Lambda_{0,\mathrm{nov}})[1]
\longrightarrow C_{3+}(L';\Lambda_{0,\mathrm{nov}})[1],\\
&\psi_*^{\mathfrak b}(\boldsymbol x)
=\mathfrak f(e^{\mathfrak b}\otimes\boldsymbol x).
\end{aligned}
\]
由于
$\Delta(e^{\mathfrak b})
=e^{\mathfrak b}\otimes e^{\mathfrak b}$，
故
\[
\widehat{\psi_*^{\mathfrak b}}(\boldsymbol x)
=\widehat{\widehat{\mathfrak f}}
 (e^{\mathfrak b}\otimes\boldsymbol x).
\]
这里的
$\widehat{\psi_*^{\mathfrak b}}\colon
BC_{1+}(L;\Lambda_{0,\mathrm{nov}})[1]
\to BC_{3+}(L';\Lambda_{0,\mathrm{nov}})[1]$
是由 $\psi_*^{\mathfrak b}$ 诱导的余代数同态，
而 $\widehat{\widehat{\mathfrak f}}$
见 (4.6.56.3)。
于是 (4.6.57) 给出
\[
\begin{aligned}
&\mathfrak q\bigl(
 e^{\psi_*\mathfrak b}
 \otimes\widehat{\psi_*^{\mathfrak b}}(\boldsymbol x)
 \bigr)\\
&\quad=
\sum_c(-1)^{\deg'\boldsymbol x_c^{(3;1)}}
\mathfrak f\Bigl(
 e^{\mathfrak b}\otimes
 \bigl(\boldsymbol x_c^{(3;1)}
 \otimes\mathfrak q(
  e^{\mathfrak b}\otimes\boldsymbol x_c^{(3;2)})
 \otimes\boldsymbol x_c^{(3;3)}\bigr)\Bigr).
\end{aligned}
\]
左边是
$\mathfrak m^{\psi_*\mathfrak b}
 (\widehat{\psi_*^{\mathfrak b}}(\boldsymbol x))$，
右边是
$\psi_*^{\mathfrak b}
 (\widehat{\mathfrak d}^{\mathfrak b}(\boldsymbol x))$。
因此 $\psi_*^{\mathfrak b}$ 给出 $A_\infty$ 同态
\[
(C_{1+}(L;\Lambda_{0,\mathrm{nov}}),
 \mathfrak m^{\mathfrak b})
\longrightarrow
(C_{3+}(L';\Lambda_{0,\mathrm{nov}}),
 \mathfrak m^{\psi_*\mathfrak b}).
\]
将系数约化到 $R$ 后，
该同态由 $\psi$ 诱导，
因而在 $\mathfrak m_1$ 上同调上为同构。
由定理 4.2.45，
它是同伦等价。

接着在 $C^2(M';\Lambda^+_{0,\mathrm{nov}})$
中找到与 $\psi_*\mathfrak b$ 同上同调的
闭上链 $\mathfrak b'$，
由定理 4.6.47，
\[
(C_{3+}(L';\Lambda_{0,\mathrm{nov}}),
  \mathfrak m^{\psi_*\mathfrak b})
\simeq
(C_{3+}(L';\Lambda_{0,\mathrm{nov}}),
  \mathfrak m^{\mathfrak b'}).
\]
最后，自然包含
\[
(C_{2+}(L';\Lambda_{0,\mathrm{nov}}),
  \mathfrak m^{\mathfrak b'})
\longrightarrow
(C_{3+}(L';\Lambda_{0,\mathrm{nov}}),
  \mathfrak m^{\mathfrak b'})
\]
也是定理 4.2.45 所保证的同伦等价。
\footnote{译者注：原书实体页 273
把这个包含两端的形变写成
$\mathfrak m^{\mathfrak b}$，
但 $\mathfrak b$ 是 $M$ 上的余链；
这里两端需使用 $M'$ 上刚选出的
$\mathfrak b'$。见疑误 A-0253。}
这样就得到了定理 4.6.52 所需的同伦等价。
同第 4.6.2 小节一样，
采用构造的“twp”版本还可证明，
它在同伦意义下不依赖于上述各种选择。
这就完成了以命题 4.6.55 为前提的
定理 4.6.52 的证明。

也可以证明定理 4.6.52 中的同伦等价
满足类似定理 4.6.44 的结论，
从而得到定理 B (B.5)。
定理 4.6.52 还有同伦幺元版本：
上述同伦等价实际上是 hu-同伦等价。
证明只须直接结合前面讨论过的技术。
